% ============================================================================
% ANHANG N: Die Suche nach magnetischen Monopolen
% ============================================================================

\chapter{Die Suche nach magnetischen Monopolen}
\label{app:monopolforschung}

\section{Einleitung}
\label{app:monopolforschung_einleitung}

Die Frage nach der Existenz magnetischer Monopole – isolierter magnetischer Ladungen – ist eine der faszinierendsten und zugleich umstrittensten Fragen der Physik. Während die Maxwell'sche Elektrodynamik das Fehlen magnetischer Monopole postuliert ($\nabla \cdot \mathbf{B} = 0$) \cite{Maxwell1873}, hat die Theorie von Paul Dirac (1931) gezeigt, dass die Existenz eines einzigen magnetischen Monopols die Quantisierung der elektrischen Ladung erklären würde \cite{Dirac1931}.

Die experimentelle Suche nach fundamentalen magnetischen Monopolen blieb jedoch über Jahrzehnte erfolglos \cite{Assis2005}. Paradoxerweise wurden in den letzten Jahren in Festkörpersystemen sogenannte \textbf{emergente magnetische Monopole} entdeckt – Quasiteilchen, die sich in ihrer Dynamik wie magnetische Monopole verhalten, aber keine fundamentalen Teilchen sind \cite{Castelnovo2008, Ladak2010, Mengotti2011}.

Dieser Anhang dokumentiert die experimentelle Monopolforschung in zwei Bereichen:
1. Die Suche nach fundamentalen Monopolen (erfolglos)
2. Die Entdeckung emergenter Monopole in Festkörpersystemen (erfolgreich)

Die Ergebnisse werden in den Kontext der Weber'schen Elektrodynamik (WED) gestellt, die fundamentale magnetische Monopole von Grund auf ausschließt (siehe Kapitel~\ref{ch:dstt_weber_kraefte} und Anhang~\ref{app:weber_metrik}).

\section{Die historische Suche nach fundamentalen Monopolen}
\label{app:monopolforschung_historisch}

Die Suche nach fundamentalen magnetischen Monopolen begann mit Diracs theoretischer Arbeit von 1931 \cite{Dirac1931} und wurde mit immer empfindlicheren Methoden fortgesetzt. Die Maxwell'sche Theorie postuliert das Fehlen magnetischer Ladungen \cite{Maxwell1873}:

\[
\nabla \cdot \mathbf{B} = 0
\]

Die Einführung magnetischer Monopole würde die Maxwell-Gleichungen symmetrisch machen \cite{Maxwell1873}:

\[
\begin{aligned}
\nabla \cdot \mathbf{E} &= \frac{\rho_e}{\varepsilon_0} \\
\nabla \cdot \mathbf{B} &= \rho_m \\
\nabla \times \mathbf{E} &= -\mu_0 \mathbf{j}_m - \frac{\partial \mathbf{B}}{\partial t} \\
\nabla \times \mathbf{B} &= \mu_0 \mathbf{j}_e + \mu_0 \varepsilon_0 \frac{\partial \mathbf{E}}{\partial t}
\end{aligned}
\]

Die Suche umfasste:
\begin{itemize}
    \item Durchmusterung kosmischer Strahlung nach Monopolen
    \item Suche in Meteoriten und Mondgestein
    \item Beschleunigerexperimente (LHC, MoEDAL)
    \item Unterirdische Detektoren (SQUID-basierte Sensoren)
\end{itemize}

Trotz intensiver Bemühungen wurde bisher kein fundamentaler magnetischer Monopol eindeutig nachgewiesen \cite{Assis2005, Assis2012}. Die Maxwell'sche Theorie mit $\nabla \cdot \mathbf{B} = 0$ bleibt damit experimentell unangefochten.

\section{Emergente magnetische Monopole in Festkörpern}
\label{app:monopolforschung_emergent}

Paradoxerweise wurden magnetische Monopole dort gefunden, wo man sie nicht erwartete – in Festkörpersystemen. Diese sogenannten \textbf{emergenten Monopole} sind Quasiteilchen, die sich in ihrer Dynamik wie magnetische Ladungen verhalten, aber keine fundamentalen Teilchen sind.

\subsection{Spin-Eis-Systeme}
\label{app:monopolforschung_spin_ice}

In sogenannten Spin-Eis-Materialien (z.B. $\mathrm{Dy_2Ti_2O_7}$) führt die geometrische Frustration der magnetischen Momente zur Entstehung von Quasiteilchen, die sich wie magnetische Monopole verhalten \cite{Castelnovo2008}. Die Bewegung dieser Monopole folgt dabei dynamischen fraktalen Trajektorien \cite{Hallen2022}.

Die experimentelle Erforschung dieser emergenten Monopole hat in den letzten Jahren bedeutende Fortschritte gemacht:

\begin{enumerate}
    \item \textbf{Monopoldynamik in Spin-Eis:} Hsu et al. (2024) untersuchten die Feld-getriebene Monopol-Dynamik in $\mathrm{Dy_2Ti_2O_7}$ und entdeckten eine dichotome Dynamik: schnelle fraktale Umordnungen gefolgt von konventioneller Monopoldiffusion \cite{Hsu2024, Blundell2024}.
    
    \item \textbf{Monopolrauschen:} Die charakteristischen Rauschspektren dieser emergenten Monopole wurden theoretisch vorhergesagt und experimentell bestätigt \cite{Dusad2019, Samarakoon2022}.
\end{enumerate}

\subsection{Künstliche Spin-Eis-Systeme}
\label{app:monopolforschung_artificial}

Neben natürlichen Spin-Eis-Materialien wurden auch künstliche Spin-Eis-Strukturen aus Nanomagneten entwickelt, in denen emergente Monopole kontrolliert werden können \cite{Ladak2010, Mengotti2011}.

Bhandari et al. (2025) demonstrierten die kontrollierte, schrittweise Bewegung emergenter Monopole in quadratischen künstlichen Spin-Eis-Strukturen \cite{Bhandari2025}:

\begin{itemize}
    \item \textbf{Kontrollierte Monopolpropagation:} Mittels "Astroid-Clocking"-Protokollen wurde die schrittweise Bewegung von Monopolen experimentell nachgewiesen.
    \item \textbf{Verschiedene Propagationsmoden:} Es wurden verschiedene Pfade der Monopolausbreitung identifiziert: String-, Zickzack- und ferromagnetische Domänen-Regime.
    \item \textbf{Anwendungen:} Diese kontrollierbaren emergenten Monopole haben Potenzial für Anwendungen in der Datenverarbeitung und im unkonventionellen Computing.
\end{itemize}

\subsection{Chiral-induzierte Monopol-ähnliche Felder}
\label{app:monopolforschung_chiral}

Eine weitere überraschende Entdeckung ist die Erzeugung monopolarer Magnetfelder durch chirale Moleküle auf superparamagnetischen Nanopartikeln \cite{Small2024}. Durch den chiral-induzierten Spin-Selektivitäts-Effekt (CISS) entstehen Nanopartikel, die im Nahfeld ein monopolarartiges Magnetfeld besitzen:

\begin{itemize}
    \item Die Richtung des Magnetfelds wird durch die Händigkeit (D- oder L-Cystein) der adsorbierten Moleküle bestimmt.
    \item Die Reichweite des monopolarartigen Feldes liegt im Nanometer-Bereich.
    \item Im Fernfeld verschwindet das Feld – es handelt sich also um einen lokalen emergenten Effekt.
\end{itemize}

\section{Die Weber'sche Perspektive auf Monopole}
\label{app:monopolforschung_weber}

Die Weber'sche Elektrodynamik (WED) nimmt zur Frage magnetischer Monopole eine klare Position ein: \textbf{Es gibt keine fundamentalen magnetischen Monopole.}

Dies folgt direkt aus der Struktur der Weber-Kraft \cite{Assis2005, Assis2012}:

\[
\vec{F} = \frac{q_1 q_2}{4\pi \varepsilon_0 r^2}
\left[
\left(
1 - \frac{\dot{r}^2}{2c^2} + \frac{r\ddot{r}}{c^2}
\right) \hat{r}
- \frac{\dot{r}^2}{2c^2} \vec{v}
\right]
\]

In der WED ist das Magnetfeld kein fundamentales Feld, sondern ein Effekt bewegter elektrischer Ladungen. Es gibt daher keine magnetischen Quellen oder Senken – die Maxwell-Gleichung $\nabla \cdot \mathbf{B} = 0$ ist keine empirische Beobachtung, sondern eine logische Konsequenz der Theorie \cite{Assis2005}:

\begin{quote}
    \emph{``There are no magnetic monopoles in Weber's electrodynamics.''} \cite{Assis2005}
\end{quote}

\subsection{Die WED und emergente Monopole}
\label{app:monopolforschung_weber_emergent}

Die Entdeckung emergenter Monopole in Spin-Eis-Systemen widerspricht der WED nicht, denn:

\begin{enumerate}
    \item \textbf{Emergente Monopole sind Quasiteilchen:} Sie sind kollektive Anregungen der magnetischen Momente in einem Festkörper, keine fundamentalen Teilchen \cite{Castelnovo2008}.
    
    \item \textbf{Das wahre Magnetfeld bleibt quellenfrei:} Das wahre Maxwell'sche Feld $\mathbf{B}$ erfüllt weiterhin $\nabla \cdot \mathbf{B} = 0$. Die emergenten Monopole erzeugen einen Fluss in einem \emph{effektiven} Magnetisierungsfeld $\mathbf{M}_{\text{eff}}$, nicht im wahren $\mathbf{B}$-Feld \cite{Bramwell2009}.
    
    \item \textbf{Die WED beschreibt die zugrundeliegende Wechselwirkung:} Die emergenten Monopole in Spin-Eis-Systemen sind eine Manifestation der zugrundeliegenden Spin-Spin-Wechselwirkungen, die durch die WED auf der fundamentalen Ebene beschrieben werden können.
\end{enumerate}

\section{Die Konsequenz für die Monopolforschung}
\label{app:monopolforschung_konsequenz}

Die experimentelle Situation lässt sich wie folgt zusammenfassen:

\begin{table}[htbp]
\centering
\begin{tabular}{|l|c|c|}
\hline
\textbf{Monopol-Typ} & \textbf{Existenz} & \textbf{Experimenteller Status} \\
\hline
Fundamentale Monopole & Nein & Bisher nicht nachgewiesen \\
Emergente Monopole (Spin-Eis) & Ja & Nachgewiesen seit 2008 \cite{Castelnovo2008} \\
Emergente Monopole (ASI) & Ja & Nachgewiesen seit 2010 \cite{Ladak2010, Mengotti2011} \\
Chirale Monopol-ähnliche Felder & Ja & Nachgewiesen seit 2024 \cite{Small2024} \\
\hline
\end{tabular}
\caption{Verschiedene Typen magnetischer Monopole und ihr experimenteller Status.}
\label{tab:monopole_status}
\end{table}

\subsection{Was bedeutet dies für die IWT/WDBT+?}
\label{app:monopolforschung_bedeutung}

Die IWT/WDBT+ sagt fundamentale magnetische Monopole nicht voraus – sie sind mit der Struktur der Theorie unvereinbar. Dies ist \textbf{kein Fehler der Theorie}, sondern eine Konsequenz ihrer konsistenten mathematischen Struktur:

\begin{enumerate}
    \item In der WED und WG ist das Magnetfeld kein fundamentales Feld (siehe Anhang~\ref{app:weber_metrik}).
    \item Die inhärente Metrik der WED/WG (Gleichung M.12) zeigt, dass die Maxwell'sche Theorie ein emergenter Grenzfall ist.
    \item Die Entdeckung emergenter Monopole in Festkörpern bestätigt indirekt die WED: Magnetische Phänomene sind auf fundamentaler Ebene durch die Wechselwirkung bewegter Ladungen (Elektronen) erklärbar.
\end{enumerate}

Die Suche nach fundamentalen Monopolen ist damit aus Sicht der WED/WG eine Sackgasse – nicht weil sie technisch schwierig wäre, sondern weil die gesuchten Teilchen in der Theorie nicht existieren. Die emergente Realität der Monopole in Festkörpersystemen zeigt jedoch, dass das Konzept der magnetischen Ladung auf einer effektiven Ebene durchaus realisierbar ist.

\section{Zusammenfassung}
\label{app:monopolforschung_zusammenfassung}

Die experimentelle Monopolforschung hat zwei Gesichter:

\begin{enumerate}
    \item \textbf{Fundamentale Monopole:} Trotz intensiver Suche nicht nachgewiesen. Die Maxwell'sche Theorie mit $\nabla \cdot \mathbf{B} = 0$ bleibt experimentell bestätigt \cite{Maxwell1873}.
    
    \item \textbf{Emergente Monopole:} In Spin-Eis-Materialien, künstlichen Spin-Eis-Strukturen und chiralen Nanopartikeln nachgewiesen \cite{Castelnovo2008, Ladak2010, Mengotti2011, Small2024}. Sie sind Quasiteilchen, die sich wie magnetische Ladungen verhalten, aber kein fundamentales Magnetfeld erzeugen.
\end{enumerate}

Die Weber'sche Elektrodynamik (WED) erklärt diese Situation konsistent:
\begin{itemize}
    \item Sie schließt fundamentale Monopole aus \cite{Assis2005}.
    \item Sie ermöglicht emergente Monopole als kollektive Effekte.
    \item Sie ist mit der Maxwell'schen Theorie kompatibel – die Maxwell-Gleichungen sind ein Grenzfall der WED \cite{Assis1994}.
\end{itemize}

Die IWT/WDBT+ geht noch einen Schritt weiter: Sie zeigt, dass die Raumzeit selbst aus Information emergiert (siehe Kapitel~\ref{ch:emergenz_raum_zeit}) und dass die Maxwell'sche Theorie ein kontinuierlicher Grenzfall einer tieferen diskreten Informationsdynamik ist (siehe Anhang~\ref{app:kontinuumslimes}). Die Monopolforschung ist damit ein Beispiel dafür, wie die moderne Physik an der Grenze zwischen fundamentaler Theorie und emergenten Phänomenen operiert – und wie die WED/WG eine konsistente Beschreibung dieser Grenze liefert.